
\documentclass[11pt]{article}
%%%%%%%%%%%%%%%%%%%%%%%%%%%%%%%%%%%%%%%%%%%%%%%%%%%%%%%%%%%%%%%%%%%%%%%%%%%%%%%%%%%%%%%%%%%%%%%%%%%%%%%%%%%%%%%%%%%%%%%%%%%%%%%%%%%%%%%%%%%%%%%%%%%%%%%%%%%%%%%%%%%%%%%%%%%%%%%%%%%%%%%%%%%%%%%%%%%%%%%%%%%%%%%%%%%%%%%%%%%%%%%%%%%%%%%%%%%%%%%%%%%%%%%%%%%%
\usepackage{xypic}
\usepackage{mathrsfs}
\usepackage{amssymb}
\usepackage{amsfonts}
\usepackage{amsmath}

\setcounter{MaxMatrixCols}{10}
%TCIDATA{OutputFilter=Latex.dll}
%TCIDATA{Version=5.50.0.2890}
%TCIDATA{<META NAME="SaveForMode" CONTENT="1">}
%TCIDATA{BibliographyScheme=Manual}
%TCIDATA{LastRevised=Monday, February 07, 2011 11:17:06}
%TCIDATA{<META NAME="GraphicsSave" CONTENT="32">}

\hyphenation{theo-ry ca-te-go-ry li-near ge-ne-ra-li-zed geo-me-try dif-fe-ren-tiable me-cha-ni-cal ma-the-ma-tics }
\textwidth146mm \textheight230mm \oddsidemargin 1.2cm
\evensidemargin 1.2cm \headsep10mm
  \topmargin-9mm
\def\ben{\begin{enumerate}}
\def\een{\end{enumerate}}
\def\dd{\displaystyle}
\def\bit{\begin{itemize}}
\def\eit{\end{itemize}}
\def\n{\noindent}
\def\mk{\medskip }
\def\bk{\bigswkip }
\def\sk{\smallskip }
\def\0{\leqno}
\def\bk{\bigskip}
\def\xym{\xymatrix}
\def\st#1#2{{\mathrel{\mathop{#2}\limits_{#1}}{}\!}}
\input xy
\xyoption{all}
\input{tcilatex}
\begin{document}

\begin{center}
\textbf{ALGEBRAIC CONSTRUCTIONS }

\textbf{IN THE CATEGORY OF VECTOR BUNDLES} \ \ \\[0pt]
\ \ %\bigskip

%\bigskip

by

%\bigskip

{CONSTANTIN M. ARCU\c{S} }
\end{center}

%\bigskip


\begin{abstract}
In this paper we present the category of generalized Lie algebroids.
Important results (a theorem of Maurer-Cartan type, theorems of Cartan
type,...) emphasize the importance and the utility of the objects of this
new category. \ \ \bigskip\newline
\textbf{2000 Mathematics Subject Classification:} 00A69, 58A15,
58B34.\bigskip\newline
\ \ \ \textbf{Keywords:} vector bundle, (generalized) Lie algebroid,
interior differential system, exterior differential calculus, exterior
differential system.
\end{abstract}
\
\
\tableofcontents
\
%\ \ \ \textit{In memory of my uncles}

%\ \ \ \textit{Prof. Dr. Gheorghe RADU and Acad. Dr. Doc. Cornelius RADU}

%\medskip
%
%\ \ \ \textit{Dedicated to} \textit{Acad. Prof. Dr. Doc. Radu MIRON at his 83%
%$^{th}$ anniversary}

%\bigskip

\ \

\section{Introduction}

The motivation for our researches was to extend the notion of Lie algebroid
using the extension $\mathbf{Mod}$-morphism associated to a $\mathbf{B}^{%
\mathbf{v}}$-morphism. Using this general framework, we get a panoramic view
over classical concepts from mathematics. [1]

We introduced the notion of \emph{interior differential system }(IDS) of a
generalized Lie algebroid, in general, and of a Lie algebroid, in
particular. We develop the exterior differential calculus for generalized
Lie algebroids and, in this general framework, we establish the structure
equations of Maurer-Cartan type and we characterize the involutivity of an
IDS in a theorem of Cartan type. Finally, using the classical notion of
\emph{exterior differential system}(see: [2,4,6,7]) (EDS) of a generalized
Lie algebroid, in general, and of a Lie algebroid, in particular, we
characterize the involutivity of an IDS in a theorem of Cartan type. In
particular, we can obtain similar results with classical results for Lie
algebroids. (see: [3,8,9])

\section{Preliminaries}

In general, if $\mathcal{C}$ is a category, then we denoted by $\left\vert
\mathcal{C}\right\vert $ the class of objects and for any $A,B{\in }%
\left\vert \mathcal{C}\right\vert $, we denote by $\mathcal{C}\left(
A,B\right) $ the set of morphisms of $A$ source and $B$ target.

Let$\mathbf{~Vect},$ $\mathbf{Liealg},~\mathbf{Mod}$\textbf{,} $\mathbf{%
Man,~B}$ and $\mathbf{B}^{\mathbf{v}}$ be the category of real vector
spaces, Lie algebras, modules, manifolds, fiber bundles and vector bundles
respectively.

We know that if $\left( E,\pi ,M\right) \in \left\vert \mathbf{B}^{\mathbf{v}%
}\right\vert ,$ then $\left( \Gamma \left( E,\pi ,M\right) ,+,\cdot \right) $
is a $\mathcal{F}\left( M\right) $-module.

In adition, if $\left( E,\pi ,M\right) \in \left\vert \mathbf{B}^{\mathbf{v}%
}\right\vert $ such that $M$ is paracompact and if $A\subseteq M$ is closed,
then for any section $u$ over $A$ it exists $\tilde{u}\in $ $\Gamma \left(
E,\pi ,M\right) $ such that $\tilde{u}_{|A}=u.$

\textbf{Note: }In the following, we consider only vector bundles with
paracompact base.

\textbf{Proposition 2.1 }If \ $\left( \varphi ,\varphi _{0}\right) \in
\mathbf{B}^{\mathbf{v}}\left( \left( E,\pi ,M\right) ,\left( E^{\prime },\pi
^{\prime },M^{\prime }\right) \right) ,$ then it exists a $\mathbf{Mod}$%
-morphism $\Gamma \left( \varphi ,\varphi _{0}\right) $ of $\Gamma \left(
E,\pi ,M\right) $ source and $\Gamma \left( E^{\prime },\pi ^{\prime
},M^{\prime }\right) $ target.

\emph{Proof. }For every $y\in \varphi _{0}\left( M\right) ,$ we fixed $%
x_{y}\in M$ such that $\varphi _{0}\left( x_{y}\right) =y.$ Then we obtain a
section $\Gamma \left( \varphi ,\varphi _{0}\right) u$ over the closed set $%
\varphi _{0}\left( M\right) $ defined by
\begin{equation*}
\begin{array}{c}
\Gamma \left( \varphi ,\varphi _{0}\right) u\left( y\right) =\varphi \left(
u_{x_{y}}\right) .%
\end{array}%
\end{equation*}

As $M^{\prime }$ is paracompact, then it results that the section $\Gamma
\left( \varphi ,\varphi _{0}\right) u$ can be regarded as a section of $%
\left( \Gamma \left( E^{\prime },\pi ^{\prime },M^{\prime }\right) +,\cdot
\right) $.\hfill

So, we obtain a $\mathbf{Mod}$-morphism%
\begin{equation*}
\begin{array}{ccc}
\Gamma \left( E,\pi ,M\right) & ^{\underrightarrow{~\ \ \Gamma \left(
\varphi ,\varphi _{0}\right) ~\ \ }} & \Gamma \left( E^{\prime },\pi
^{\prime },M^{\prime }\right) \\
u & \longmapsto & \Gamma \left( \varphi ,\varphi _{0}\right) u%
\end{array}%
\end{equation*}%
defined by
\begin{equation*}
\begin{array}{c}
\Gamma \left( \varphi ,\varphi _{0}\right) u\left( y\right) =\varphi \left(
u_{x_{y}}\right) ,%
\end{array}%
\end{equation*}%
for any $y\in \varphi _{0}\left( M\right) .$\hfill \emph{q.e.d.}\medskip

\textbf{Definition 2.1 }A $\mathbf{Mod}$-morphism given by the previous
proposition is called \emph{the extension }$\mathbf{Mod}$\emph{-morphism
associated to the }$\mathbf{B}^{\mathbf{v}}$\emph{-morphism }$\left( \varphi
,\varphi _{0}\right) .$

\textit{Remark 2.1 }The construction of the extension\emph{\ }$\mathbf{Mod}$%
-morphism associated to a $\mathbf{B}^{\mathbf{v}}$-morphism $\left( \varphi
,\varphi _{0}\right) $ is not unique, but any two extension $\mathbf{Mod}$%
-morphisms associated to a $\mathbf{B}^{\mathbf{v}}$-morphism $\left(
\varphi ,\varphi _{0}\right) $ has the same properties.

\textbf{Example 2.1 }If \ $\left( \varphi ,\varphi _{0}\right) \in \mathbf{B}%
^{\mathbf{v}}\left( \left( E,\pi ,M\right) ,\left( E^{\prime },\pi ^{\prime
},M^{\prime }\right) \right) $ such that $\varphi _{0}\in Diff\left(
M,M^{\prime }\right) ,$ then we obtain the unique $\mathbf{Mod}$-morphism%
\begin{equation*}
\begin{array}{ccc}
\Gamma \left( E,\pi ,M\right) & ^{\underrightarrow{~\ \ \Gamma \left(
\varphi ,\varphi _{0}\right) ~\ \ }} & \Gamma \left( E^{\prime },\pi
^{\prime },M^{\prime }\right) \\
u & \longmapsto & \Gamma \left( \varphi ,\varphi _{0}\right) u%
\end{array}%
\end{equation*}%
defined by
\begin{equation*}
\begin{array}{c}
\left( \Gamma \left( \varphi ,\varphi _{0}\right) u\right) \left( x^{\prime
}\right) =\varphi \left( u_{\varphi _{0}^{-1}\left( x^{\prime }\right)
}\right) .%
\end{array}%
\end{equation*}

$\Gamma \left( \varphi ,\varphi _{0}\right) $ is called \emph{the }$\mathbf{%
Mod}$\emph{-morphism associated to the }$\mathbf{B}^{\mathbf{v}}$\emph{%
-morphism }$\left( \varphi ,\varphi _{0}\right) .$

\section{The category of generalized Lie algebroids}

We know that a Lie algebroid is a vector bundle $\left( F,\nu ,N\right) \in
\left\vert \mathbf{B}^{\mathbf{v}}\right\vert $ such that there exists
\begin{equation*}
\begin{array}{c}
\left( \rho ,Id_{N}\right) \in \mathbf{B}^{\mathbf{v}}\left( \left( F,\nu
,N\right) ,\left( TN,\tau _{N},N\right) \right)%
\end{array}%
\end{equation*}%
and an operation
\begin{equation*}
\begin{array}{ccc}
\Gamma \left( F,\nu ,N\right) \times \Gamma \left( F,\nu ,N\right) & ^{%
\underrightarrow{\,\left[ ,\right] _{F}\,}} & \Gamma \left( F,\nu ,N\right)
\\
\left( u,v\right) & \longmapsto & \left[ u,v\right] _{F}%
\end{array}%
\end{equation*}%
with the following properties:

\begin{itemize}
\item[$LA_{1}$.] the equality holds good
\begin{equation*}
\begin{array}{c}
\left[ u,f\cdot v\right] _{F}=f\left[ u,v\right] _{F}+\Gamma \left( \rho
,Id_{N}\right) \left( u\right) f\cdot v,%
\end{array}%
\end{equation*}%
for all $u,v\in \Gamma \left( F,\nu ,N\right) $ and $f\in \mathcal{F}\left(
N\right) ,$

\item[$LA_{2}$.] the $4$-tuple $\left( \Gamma \left( F,\nu ,N\right)
,+,\cdot ,\left[ ,\right] _{F}\right) $ is a Lie $\mathcal{F}\left( N\right)
$-algebra$,$

\item[$LA_{3}$.] the $\mathbf{Mod}$-morphism $\Gamma \left( \rho
,Id_{N}\right) $ is a $\mathbf{LieAlg}$-morphism of $\left( \Gamma \left(
F,\nu ,N\right) ,+,\cdot ,\left[ ,\right] _{F}\right) $ source and $\left(
\Gamma \left( TN,\tau _{N},N\right) ,+,\cdot ,\left[ ,\right] _{TN}\right) $
target.
\end{itemize}

Obviously, in the definition of the Lie algebroid we use essential the $%
\mathbf{Mod}$-morphism $\Gamma \left( \rho ,Id_{N}\right) $ associated to
the $\mathbf{B}^{\mathbf{v}}$-morphism $\left( \rho ,Id_{N}\right) .$

So, we are interested to finding the answer to the following question:

\begin{itemize}
\item \emph{Could we to extend the notion of Lie algebroid using the
extension }$\mathbf{Mod}$\emph{-morphism associated to a }$\mathbf{B}^{%
\mathbf{v}}$\emph{-morphism?}
\end{itemize}

\textbf{Definition 3.1 }Let $M,N\in \left\vert \mathbf{Man}\right\vert $ and
$h\in \mathbf{Man}\left( M,N\right) $ be surjective.

If $\left( F,\nu ,N\right) \in \left\vert \mathbf{B}^{\mathbf{v}}\right\vert
$ such that there exists
\begin{equation*}
\begin{array}{c}
\left( \rho ,\eta \right) \in \mathbf{B}^{\mathbf{v}}\left( \left( F,\nu
,N\right) ,\left( TM,\tau _{M},M\right) \right)%
\end{array}%
\end{equation*}%
and an operation
\begin{equation*}
\begin{array}{ccc}
\Gamma \left( F,\nu ,N\right) \times \Gamma \left( F,\nu ,N\right) & ^{%
\underrightarrow{\left[ ,\right] _{F,h}}} & \Gamma \left( F,\nu ,N\right) \\
\left( u,v\right) & \longmapsto & \left[ u,v\right] _{F,h}%
\end{array}%
\end{equation*}%
with the following properties:\bigskip

\noindent $\qquad GLA_{1}$. the equality holds good
\begin{equation*}
\begin{array}{c}
\left[ u,f\cdot v\right] _{F,h}=f\left[ u,v\right] _{F,h}+\Gamma \left(
Th\circ \rho ,h\circ \eta \right) \left( u\right) f\cdot v,%
\end{array}%
\end{equation*}%
\qquad \quad\ \ for all $u,v\in \Gamma \left( F,\nu ,N\right) $ and $f\in
\mathcal{F}\left( N\right) .$

\medskip $GLA_{2}$. the $4$-tuple $\left( \Gamma \left( F,\nu ,N\right)
,+,\cdot ,\left[ ,\right] _{F,h}\right) $ is a Lie $\mathcal{F}\left(
N\right) $-algebra,

$GLA_{3}$. the $\mathbf{Mod}$-morphism $\Gamma \left( Th\circ \rho ,h\circ
\eta \right) $ is a $\mathbf{LieAlg}$-morphism of $\left( \Gamma \left(
F,\nu ,N\right) ,+,\cdot ,\left[ ,\right] _{F,h}\right) $ source and $\left(
\Gamma \left( TN,\tau _{N},N\right) ,+,\cdot ,\left[ ,\right] _{TN}\right) $
target, \medskip \noindent then we will say that \emph{the triple }%
\begin{equation*}
\begin{array}{c}
\left( \left( F,\nu ,N\right) ,\left[ ,\right] _{F,h},\left( \rho ,\eta
\right) \right)%
\end{array}%
\leqno(3.1)
\end{equation*}%
\emph{is a generalized Lie algebroid. }The couple $\left( \left[ ,\right]
_{F,h},\left( \rho ,\eta \right) \right) $will be called \emph{generalized
Lie algebroid structure.}

\textbf{Definition 3.2 }We define the set of morphisms of
\begin{equation*}
\left( \left( F,\nu ,N\right) ,\left[ ,\right] _{F,h},\left( \rho ,\eta
\right) \right)
\end{equation*}%
source and
\begin{equation*}
\left( \left( F^{\prime },\nu ^{\prime },N^{\prime }\right) ,\left[ ,\right]
_{F^{\prime },h^{\prime }},\left( \rho ^{\prime },\eta ^{\prime }\right)
\right)
\end{equation*}%
target as being the set
\begin{equation*}
\begin{array}{c}
\left\{ \left( \varphi ,\varphi _{0}\right) \in \mathbf{B}^{\mathbf{v}%
}\left( \left( F,\nu ,N\right) ,\left( F^{\prime },\nu ^{\prime },N^{\prime
}\right) \right) \right\}%
\end{array}%
\end{equation*}%
such that the $\mathbf{Mod}$-morphism $\Gamma \left( \varphi ,\varphi
_{0}\right) $ is a $\mathbf{LieAlg}$-morphism of
\begin{equation*}
\left( \Gamma \left( F,\nu ,N\right) ,+,\cdot ,\left[ ,\right] _{F,h}\right)
\end{equation*}%
source and
\begin{equation*}
\left( \Gamma \left( F^{\prime },\nu ^{\prime },N^{\prime }\right) ,+,\cdot
, \left[ ,\right] _{F^{\prime },h^{\prime }}\right)
\end{equation*}%
target.

We remark that we can discuss about \emph{the category }$\mathbf{GLA}$\emph{%
\ of generalized Lie algebroids.}

Let $\left( \left( F,\nu ,N\right) ,\left[ ,\right] _{F,h},\left( \rho ,\eta
\right) \right) $ be a generalized Lie algebroid.

\begin{itemize}
\item Locally, for any $\alpha ,\beta \in \overline{1,p},$ we set $\left[
t_{\alpha },t_{\beta }\right] _{F}\overset{put}{=}L_{\alpha \beta }^{\gamma
}t_{\gamma }.$ We easily obtain that $L_{\alpha \beta }^{\gamma }=-L_{\beta
\alpha }^{\gamma },~$for any $\alpha ,\beta ,\gamma \in \overline{1,p}.$
\end{itemize}

The real local functions $\left\{ L_{\alpha \beta }^{\gamma },~\alpha ,\beta
,\gamma \in \overline{1,p}\right\} $ will be called the \emph{structure
functions of the generalized Lie algebroid }$\left( \left( F,\nu ,N\right) ,%
\left[ ,\right] _{F,h},\left( \rho ,\eta \right) \right) .$

\begin{itemize}
\item We assume that $\left( F,\nu ,N\right) $ is a vector bundle with type
fibre the real vector space $\left( \mathbb{R}^{p},+,\cdot \right) $ and
structure group a Lie subgroup of $\left( \mathbf{GL}\left( p,\mathbb{R}%
\right) ,\cdot \right) .$
\end{itemize}

We take $\left( x^{i},y^{i}\right) $ as canonical local coordinates on $%
\left( TM,\tau _{M},M\right) ,$ where $i\in \overline{1,m}.$

Consider $\left( x^{i},y^{i}\right) \longrightarrow \left( x^{i%
%TCIMACRO{\U{b4}}%
%BeginExpansion
{\acute{}}%
%EndExpansion
}\left( x^{i}\right) ,y^{i%
%TCIMACRO{\U{b4}}%
%BeginExpansion
{\acute{}}%
%EndExpansion
}\left( x^{i},y^{i}\right) \right) $ a change of coordinates on $\left(
TM,\tau _{M},M\right) $. Then the coordinates $y^{i}$ change to $y^{i%
%TCIMACRO{\U{b4}}%
%BeginExpansion
{\acute{}}%
%EndExpansion
}$ by the rule:
\begin{equation*}
\begin{array}{c}
y^{i%
%TCIMACRO{\U{b4}}%
%BeginExpansion
{\acute{}}%
%EndExpansion
}=\frac{\partial x^{i%
%TCIMACRO{\U{b4}}%
%BeginExpansion
{\acute{}}%
%EndExpansion
}}{\partial x^{i}}y^{i}.%
\end{array}%
\leqno(3.2)
\end{equation*}

We take $(\varkappa ^{\tilde{\imath}},z^{\alpha })$ as canonical local
coordinates on $(F,\nu ,N),$ where $\tilde{\imath}{\in }\overline{1,n}$, $%
\alpha \in \overline{1,p}.$

Consider $\left( \varkappa ^{\tilde{\imath}},z^{\alpha }\right)
\longrightarrow \left( \varkappa ^{\tilde{\imath}%
%TCIMACRO{\U{b4}}%
%BeginExpansion
{\acute{}}%
%EndExpansion
},z^{\alpha
%TCIMACRO{\U{b4}}%
%BeginExpansion
{\acute{}}%
%EndExpansion
}\right) $ a change of coordinates on $\left( F,\nu ,N\right) $. Then the
coordinates $z^{\alpha }$ change to $z^{\alpha
%TCIMACRO{\U{b4}}%
%BeginExpansion
{\acute{}}%
%EndExpansion
}$ by the rule:
\begin{equation*}
\begin{array}{c}
z^{\alpha
%TCIMACRO{\U{b4}}%
%BeginExpansion
{\acute{}}%
%EndExpansion
}=\Lambda _{\alpha }^{\alpha
%TCIMACRO{\U{b4}}%
%BeginExpansion
{\acute{}}%
%EndExpansion
}z^{\alpha }.%
\end{array}%
\leqno(3.3)
\end{equation*}

\begin{itemize}
\item We assume that $\left( \theta ,\mu \right) \overset{put}{=}\left(
Th\circ \rho ,h\circ \eta \right) $. If $z^{\alpha }t_{\alpha }\in \Gamma
\left( F,\nu ,N\right) $ is arbitrary, then
\begin{equation*}
\begin{array}{c}
\Gamma \left( Th\circ \rho ,h\circ \eta \right) \left( z^{\alpha }t_{\alpha
}\right) f\left( h\circ \eta \left( \varkappa \right) \right) =\vspace*{1mm}
\\
\displaystyle=\left( \theta _{\alpha }^{\tilde{\imath}}z^{\alpha }\frac{%
\partial f}{\partial \varkappa ^{\tilde{\imath}}}\right) \left( h\circ \eta
\left( \varkappa \right) \right) =\left( \left( \rho _{\alpha }^{i}\circ
h\right) \left( z^{\alpha }\circ h\right) \frac{\partial f\circ h}{\partial
x^{i}}\right) \left( \eta \left( \varkappa \right) \right) ,%
\end{array}%
\leqno(3.4)
\end{equation*}%
for any $f\in \mathcal{F}\left( N\right) $ and $\varkappa \in N.$
\end{itemize}

The coefficients $\rho _{\alpha }^{i}$ respectively $\theta _{\alpha }^{%
\tilde{\imath}}$ change to $\rho _{\alpha
%TCIMACRO{\U{b4}}%
%BeginExpansion
{\acute{}}%
%EndExpansion
}^{i%
%TCIMACRO{\U{b4}}%
%BeginExpansion
{\acute{}}%
%EndExpansion
}$ respectively $\theta _{\alpha
%TCIMACRO{\U{b4}}%
%BeginExpansion
{\acute{}}%
%EndExpansion
}^{\tilde{\imath}%
%TCIMACRO{\U{b4}}%
%BeginExpansion
{\acute{}}%
%EndExpansion
}$ by the rule:
\begin{equation*}
\begin{array}{c}
\rho _{\alpha
%TCIMACRO{\U{b4}}%
%BeginExpansion
{\acute{}}%
%EndExpansion
}^{i%
%TCIMACRO{\U{b4}}%
%BeginExpansion
{\acute{}}%
%EndExpansion
}=\Lambda _{\alpha
%TCIMACRO{\U{b4}}%
%BeginExpansion
{\acute{}}%
%EndExpansion
}^{\alpha }\rho _{\alpha }^{i}\displaystyle\frac{\partial x^{i%
%TCIMACRO{\U{b4}}%
%BeginExpansion
{\acute{}}%
%EndExpansion
}}{\partial x^{i}},%
\end{array}%
\leqno(3.5)
\end{equation*}%
respectively%
\begin{equation*}
\begin{array}{c}
\theta _{\alpha
%TCIMACRO{\U{b4}}%
%BeginExpansion
{\acute{}}%
%EndExpansion
}^{\tilde{\imath}%
%TCIMACRO{\U{b4}}%
%BeginExpansion
{\acute{}}%
%EndExpansion
}=\Lambda _{\alpha
%TCIMACRO{\U{b4}}%
%BeginExpansion
{\acute{}}%
%EndExpansion
}^{\alpha }\theta _{\alpha }^{\tilde{\imath}}\displaystyle\frac{\partial
\varkappa ^{\tilde{\imath}%
%TCIMACRO{\U{b4}}%
%BeginExpansion
{\acute{}}%
%EndExpansion
}}{\partial \varkappa ^{\tilde{\imath}}},%
\end{array}%
\leqno(3.6)
\end{equation*}%
where $\left\Vert \Lambda _{\alpha
%TCIMACRO{\U{b4}}%
%BeginExpansion
{\acute{}}%
%EndExpansion
}^{\alpha }\right\Vert =\left\Vert \Lambda _{\alpha }^{\alpha
%TCIMACRO{\U{b4}}%
%BeginExpansion
{\acute{}}%
%EndExpansion
}\right\Vert ^{-1}.$

\textit{Remark 3.1 }\emph{The following equalities hold good:}%
\begin{equation*}
\begin{array}{c}
\displaystyle\rho _{\alpha }^{i}\circ h\frac{\partial f\circ h}{\partial
x^{i}}=\left( \theta _{\alpha }^{\tilde{\imath}}\frac{\partial f}{\partial
\varkappa ^{\tilde{\imath}}}\right) \circ h,\forall f\in \mathcal{F}\left(
N\right) .%
\end{array}%
\leqno(3.7)
\end{equation*}%
\emph{and }%
\begin{equation*}
\begin{array}{c}
\displaystyle\left( L_{\alpha \beta }^{\gamma }\circ h\right) \left( \rho
_{\gamma }^{k}\circ h\right) =\left( \rho _{\alpha }^{i}\circ h\right) \frac{%
\partial \left( \rho _{\beta }^{k}\circ h\right) }{\partial x^{i}}-\left(
\rho _{\beta }^{j}\circ h\right) \frac{\partial \left( \rho _{\alpha
}^{k}\circ h\right) }{\partial x^{j}}.%
\end{array}%
\leqno(3.8)
\end{equation*}

In the next we build some examples of objects of the category $\mathbf{GLA.}$

\textbf{Theorem 3.1} \emph{Let }$M\in \left\vert \mathbf{Man}_{m}\right\vert
$\emph{\ and }$g,h\in Iso_{\mathbf{Man}}\left( M\right) $\emph{\ be. Using
the tangent }$\mathbf{B}^{\mathbf{v}}$\emph{-morphism }$\left( Tg,g\right) $%
\emph{\ and the operation }%
\begin{equation*}
\begin{array}{ccc}
\Gamma \left( TM,\tau _{M},M\right) \times \Gamma \left( TM,\tau
_{M},M\right) & ^{\underrightarrow{~\ \ \left[ ,\right] _{TM,h}~\ \ }} &
\Gamma \left( TM,\tau _{M},M\right) \\
\left( u,v\right) & \longmapsto & \ \left[ u,v\right] _{TM,h}%
\end{array}%
\end{equation*}%
\emph{where }%
\begin{equation*}
\left[ u,v\right] _{TM,h}=\Gamma \left( T\left( h\circ g\right) ^{-1},\left(
h\circ g\right) ^{-1}\right) \left( \left[ \Gamma \left( T\left( h\circ
g\right) ,h\circ g\right) u,\Gamma \left( T\left( h\circ g\right) ,h\circ
g\right) v\right] _{TM}\right) ,
\end{equation*}%
\emph{for any }$u,v\in \Gamma \left( TM,\tau _{M},M\right) $\emph{, we
obtain that}%
\begin{equation*}
\begin{array}{c}
\left( \left( TM,\tau _{M},M\right) ,\left( Tg,g\right) ,\left[ ,\right]
_{TM,h}\right) \in \left\vert \mathbf{GLA}\right\vert .%
\end{array}%
\end{equation*}

For any $\mathbf{Man}$-isomorphisms $g$ and $h$ we obtain new and
interesting generalized Lie algebroid structures for the tangent vector
bundle $\left( TM,\tau _{M},M\right) .$ For any base $\left\{ t_{\alpha
},~\alpha \in \overline{1,m}\right\} $ of the module of sections $\left(
\Gamma \left( TM,\tau _{M},M\right) ,+,\cdot \right) $ we obtain the
structure functions%
\begin{equation*}
\begin{array}{c}
L_{\alpha \beta }^{\gamma }=\left( \theta _{\alpha }^{i}\frac{\partial
\theta _{\beta }^{j}}{\partial x^{i}}-\theta _{\beta }^{i}\frac{\partial
\theta _{\alpha }^{j}}{\partial x^{i}}\right) \tilde{\theta}_{j}^{\gamma
},~\alpha ,\beta ,\gamma \in \overline{1,m}%
\end{array}%
\end{equation*}%
where $\theta _{\alpha }^{i},~i,\alpha \in \overline{1,m}$ are real local
functions such that
\begin{equation*}
\begin{array}{c}
\Gamma \left( T\left( h\circ g\right) ,h\circ g\right) \left( t_{\alpha
}\right) =\theta _{\alpha }^{i}\frac{\partial }{\partial x^{i}}%
\end{array}%
\end{equation*}%
and $\tilde{\theta}_{j}^{\gamma },~i,\gamma \in \overline{1,m}$ are real
local functions such that
\begin{equation*}
\begin{array}{c}
\Gamma \left( T\left( h\circ g\right) ^{-1},\left( h\circ g\right)
^{-1}\right) \left( \frac{\partial }{\partial x^{j}}\right) =\tilde{\theta}%
_{j}^{\gamma }t_{\gamma }.%
\end{array}%
\end{equation*}

In particular, using arbitrary basis for the module of sections and
arbitrary isometries (symmetries, translations, rotations,...) for the
Euclidean $3$-dimensional space $\Sigma ,$ we obtain a lot of generalized
Lie algebroid structures for the tangent vector bundle $\left( T\Sigma ,\tau
_{\Sigma },\Sigma \right) $.

We assume that $\left( \left( F,\nu ,N\right) ,\left[ ,\right] _{F},\left(
\rho ,Id_{N}\right) \right) $ is a Lie algebroid and let $h\in \mathbf{Man}%
\left( N,N\right) $ be a surjective application. Let $\mathcal{AF}_{F}$ be a
vector fibred $\left( n+p\right) $-atlas for the vector bundle $\left( F,\nu
,N\right) $ and let $\mathcal{AF}_{TN}$ be a vector fibred $\left(
n+n\right) $-atlas for the vector bundle $\left( TN,\tau _{N},N\right) $.

If $\left( U,\xi _{U}\right) \in \mathcal{AF}_{TN}$ and $\left(
V,s_{V}\right) \in \mathcal{AF}_{F}$ such that $U\cap h^{-1}\left( V\right)
\neq \phi $, then we define the application%
\begin{equation*}
\begin{array}{ccc}
\tau _{N}^{-1}(U{\cap }h^{-1}(V))) & {}^{\underrightarrow{\bar{\xi}_{U{\cap }%
h^{-1}(V)}}} & \left( U{\cap }h^{-1}(V)\right) {\times }\mathbb{R}^{n} \\
\left( \varkappa ,u\left( \varkappa \right) \right) & \longmapsto & \left(
\varkappa ,\xi _{U,\varkappa }^{-1}u\left( \varkappa \right) \right) .%
\end{array}%
\end{equation*}

\textbf{Proposition 3.1 }\emph{The set}%
\begin{equation*}
\begin{array}{c}
\overline{\mathcal{AF}}_{TN}\overset{put}{=}\underset{U\cap h^{-1}\left(
V\right) \neq \phi }{\underset{\left( U,\xi _{U}\right) \in \mathcal{AF}%
_{TN},~\left( V,s_{V}\right) \in \mathcal{AF}_{F}}{\tbigcup }}\left\{ \left(
U\cap h^{-1}\left( V\right) ,\bar{\xi}_{U{\cap }h^{-1}(V)}\right) \right\}%
\end{array}%
\end{equation*}%
\emph{is a vector fibred }$n+n$\emph{-atlas for the vector bundle }$\left(
TN,\tau _{N},N\right) .$

\emph{If }$X=X^{\tilde{\imath}}\frac{\partial }{\partial \varkappa ^{\tilde{%
\imath}}}\in \Gamma \left( TN,\tau _{N},N\right) ,$ \emph{then\ we obtain
the section }%
\begin{equation*}
\begin{array}{c}
\bar{X}=\bar{X}^{\tilde{\imath}}\circ h\frac{\partial }{\partial \bar{%
\varkappa}^{\tilde{\imath}}}\in \Gamma \left( TN,\tau _{N},N\right) ,%
\end{array}%
\end{equation*}%
\emph{such that }$\bar{X}\left( \bar{\varkappa}\right) =X\left( h\left( \bar{%
\varkappa}\right) \right) ,$ \emph{for any }$\bar{\varkappa}\in U\cap
h^{-1}\left( V\right) .$

The set $\left\{ \frac{\partial }{\partial \bar{\varkappa}^{\tilde{\imath}}}%
,~\tilde{\imath}\in \overline{1,n}\right\} $ is a base for the $\mathcal{F}%
\left( N\right) $-module $\left( \Gamma \left( TN,\tau _{N},N\right)
,+,\cdot \right) .$

\textbf{Theorem 3.2} \emph{If we consider the operation }%
\begin{equation*}
\begin{array}{ccc}
\Gamma \left( F,\nu ,N\right) \times \Gamma \left( F,\nu ,N\right) & ^{%
\underrightarrow{~\ \ \left[ ,\right] _{F,h}~\ \ }} & \Gamma \left( F,\nu
,N\right)%
\end{array}%
\end{equation*}%
\emph{defined by}%
\begin{equation*}
\begin{array}{ll}
\left[ t_{\alpha },t_{\beta }\right] _{F,h} & =\left( L_{\alpha \beta
}^{\gamma }\circ h\right) t_{\gamma },\vspace*{1mm} \\
\left[ t_{\alpha },ft_{\beta }\right] _{F,h} & =f\left( L_{\alpha \beta
}^{\gamma }\circ h\right) t_{\gamma }+\rho _{\alpha }^{\tilde{\imath}}\circ h%
\displaystyle\frac{\partial f}{\partial \bar{\varkappa}^{\tilde{\imath}}}%
t_{\beta },\vspace*{1mm} \\
\left[ ft_{\alpha },t_{\beta }\right] _{F,h} & =-\left[ t_{\beta
},ft_{\alpha }\right] _{F,h},~%
\end{array}%
\end{equation*}%
\emph{for any }$f\in \mathcal{F}\left( N\right) ,$ \emph{then }$\left(
\left( F,\nu ,N\right) ,\left[ ,\right] _{F,h},\left( \rho ,Id_{N}\right)
\right) \in \left\vert \mathbf{GLA}\right\vert .$

The generalized Lie algebroid
\begin{equation*}
\begin{array}{c}
\left( \left( F,\nu ,N\right) ,\left[ ,\right] _{F,h},\left( \rho
,Id_{N}\right) \right)%
\end{array}%
\end{equation*}%
given by the previous theorem, will be called \emph{the generalized Lie
algebroid associated to the Lie algebroid }$\left( \left( F,\nu ,N\right) ,%
\left[ ,\right] _{F},\left( \rho ,Id_{N}\right) \right) $ \emph{and to the
surjective application }$h\in \mathbf{Man}\left( N,N\right) .$

In particular, if $h=Id_{N},$ then the generalized Lie algebroid
\begin{equation*}
\begin{array}{c}
\left( \left( F,\nu ,N\right) ,\left[ ,\right] _{F,Id_{N}},\left( \rho
,Id_{N}\right) \right)%
\end{array}%
\end{equation*}%
will be called \emph{the generalized Lie algebroid associated to the Lie
algebroid}%
\begin{equation*}
\begin{array}{c}
\left( \left( F,\nu ,N\right) ,\left[ ,\right] _{F},\left( \rho
,Id_{N}\right) \right) .%
\end{array}%
\end{equation*}

\subsection{The pull-back Lie algebroid of a generalized Lie algebroid}

Let $\left( \left( F,\nu ,N\right) ,\left[ ,\right] _{F,h},\left( \rho ,\eta
\right) \right) $ be a generalized Lie algebroid.

Let $\mathcal{AF}_{F}$ be a vector fibred $\left( n+p\right) $-atlas for the
vector bundle $\left( F,\nu ,N\right) $ and let $\mathcal{AF}_{TM}$ be a
vector fibred $\left( m+m\right) $-atlas for the vector bundle $\left(
TM,\tau _{M},M\right) $.

Let $\left( h^{\ast }F,h^{\ast }\nu ,M\right) $ be the pull-back vector
bundle through $h.$

If $\left( U,\xi _{U}\right) \in \mathcal{AF}_{TM}$ and $\left(
V,s_{V}\right) \in \mathcal{AF}_{F}$ such that $U\cap h^{-1}\left( V\right)
\neq \phi $, then we define the application%
\begin{equation*}
\begin{array}{ccc}
h^{\ast }\nu ^{-1}(U{\cap }h^{-1}(V))) & {}^{\underrightarrow{\bar{s}_{U{%
\cap }h^{-1}(V)}}} & \left( U{\cap }h^{-1}(V)\right) {\times }\mathbb{R}^{p}
\\
\left( \varkappa ,z\left( h\left( \varkappa \right) \right) \right) &
\longmapsto & \left( \varkappa ,t_{V,h\left( \varkappa \right) }^{-1}z\left(
h\left( \varkappa \right) \right) \right).%
\end{array}%
\end{equation*}

\textbf{Proposition 3.1.1 }\emph{The set}%
\begin{equation*}
\begin{array}{c}
\overline{\mathcal{AF}}_{F}\overset{put}{=}\underset{U\cap h^{-1}\left(
V\right) \neq \phi }{\underset{\left( U,\xi _{U}\right) \in \mathcal{AF}%
_{TM},~\left( V,s_{V}\right) \in \mathcal{AF}_{F}}{\tbigcup }}\left\{ \left(
U\cap h^{-1}\left( V\right) ,\bar{s}_{U{\cap }h^{-1}(V)}\right) \right\}%
\end{array}%
\end{equation*}%
\emph{is a vector fibred }$m+p$\emph{-atlas for the vector bundle }$\left(
h^{\ast }F,h^{\ast }\nu ,M\right) .$

\emph{If }$z=z^{\alpha }t_{\alpha }\in \Gamma \left( F,\nu ,N\right) ,$
\emph{then\ we obtain the section }%
\begin{equation*}
\begin{array}{c}
Z=\left( z^{\alpha }\circ h\right) T_{\alpha }\in \Gamma \left( h^{\ast
}F,h^{\ast }\nu ,M\right)
\end{array}%
\end{equation*}%
\emph{such that }$Z\left( x\right) =z\left( h\left( x\right) \right) ,$
\emph{for any }$x\in U\cap h^{-1}\left( V\right) .$\bigskip

\textbf{Theorem 3.1.1 }\emph{Let} $\Big({\overset{h^{\ast }F}{\rho }},Id_{M}%
\Big)$ \emph{be the }$\mathbf{B}^{\mathbf{v}}$\emph{-morphism of }$\left(
h^{\ast }F,h^{\ast }\nu ,M\right) $\ \emph{source and} $\left( TM,\tau
_{M},M\right) $\ \emph{target, where}%
\begin{equation*}
\begin{array}{rcl}
h^{\ast }F & ^{\underrightarrow{\overset{h^{\ast }F}{\rho }}} & TM \\
\displaystyle Z^{\alpha }T_{\alpha }\left( x\right) & \longmapsto & %
\displaystyle\left( Z^{\alpha }\cdot \rho _{\alpha }^{i}\circ h\right) \frac{%
\partial }{\partial x^{i}}\left( x\right)%
\end{array}%
\leqno(3.1.1)
\end{equation*}

\emph{Using the operation}
\begin{equation*}
\begin{array}{ccc}
\Gamma \left( h^{\ast }F,h^{\ast }\nu ,M\right) \times \Gamma \left( h^{\ast
}F,h^{\ast }\nu ,M\right) & ^{\underrightarrow{~\ \ \left[ ,\right]
_{h^{\ast }F}~\ \ }} & \Gamma \left( h^{\ast }F,h^{\ast }\nu ,M\right)%
\end{array}%
\end{equation*}%
\emph{defined by}%
\begin{equation*}
\begin{array}{ll}
\left[ T_{\alpha },T_{\beta }\right] _{h^{\ast }F} & =\left( L_{\alpha \beta
}^{\gamma }\circ h\right) T_{\gamma },\vspace*{1mm} \\
\left[ T_{\alpha },fT_{\beta }\right] _{h^{\ast }F} & \displaystyle=f\left(
L_{\alpha \beta }^{\gamma }\circ h\right) T_{\gamma }+\left( \rho _{\alpha
}^{i}\circ h\right) \frac{\partial f}{\partial x^{i}}T_{\beta },\vspace*{1mm}
\\
\left[ fT_{\alpha },T_{\beta }\right] _{h^{\ast }F} & =-\left[ T_{\beta
},fT_{\alpha }\right] _{h^{\ast }F},%
\end{array}%
\leqno(3.1.2)
\end{equation*}%
\emph{for any} $f\in \mathcal{F}\left( M\right) ,$ \emph{it results that}
\begin{equation*}
\begin{array}{c}
\left( \left( h^{\ast }F,h^{\ast }\nu ,M\right) ,\left[ ,\right] _{h^{\ast
}F},\left( \overset{h^{\ast }F}{\rho },Id_{M}\right) \right)%
\end{array}%
\end{equation*}%
ia a Lie algebroid which is called \emph{the pull-back Lie algebroid of the
generalized Lie algebroid }$\left( \left( F,\nu ,N\right) ,\left[ ,\right]
_{F,h},\left( \rho ,\eta \right) \right) .$

\subsection{Interior Differential Systems}

Let $\left( \left( h^{\ast }F,h^{\ast }\nu ,M\right) ,\left[ ,\right]
_{h^{\ast }F},\left( \overset{h^{\ast }F}{\rho },Id_{M}\right) \right) $ be
the pull-back Lie algebroid of the generalized Lie algebroid $\left( \left(
F,\nu ,N\right) ,\left[ ,\right] _{F,h},\left( \rho ,\eta \right) \right) $.

\textbf{Definition 3.2.1 }Any vector subbundle $\left( E,\pi ,M\right) $ of
the vector bundle $\left( h^{\ast }F,h^{\ast }\nu ,M\right) $ will be called
\emph{interior differential system (IDS) of the generalized Lie algebroid }%
\begin{equation*}
\left( \left( F,\nu ,N\right) ,\left[ ,\right] _{F,h},\left( \rho ,\eta
\right) \right) .
\end{equation*}

In particular, if $h=Id_{N}=\eta $, then any vector subbundle $\left( E,\pi
,N\right) $ of the vector bundle $\left( F,\nu ,N\right) $ will be called
\emph{interior differential system of the Lie algebroid }%
\begin{equation*}
\left( \left( F,\nu ,N\right) ,\left[ ,\right] _{F},\left( \rho
,Id_{N}\right) \right) .
\end{equation*}

\textit{Remark 3.2.1 }If $\left( E,\pi ,M\right) $ is an IDS of the
generalized Lie algebroid
\begin{equation*}
\left( \left( F,\nu ,N\right) ,\left[ ,\right] _{F,h},\left( \rho ,\eta
\right) \right) ,
\end{equation*}%
then we obtain a vector subbundle $\left( E^{0},\pi ^{0},M\right) $ of the
vector bundle $\left( \overset{\ast }{h^{\ast }F},\overset{\ast }{h^{\ast
}\nu },M\right) $ such that
\begin{equation*}
\Gamma \left( E^{0},\pi ^{0},M\right) \overset{put}{=}\left\{ \Omega \in
\Gamma \left( \overset{\ast }{h^{\ast }F},\overset{\ast }{h^{\ast }\nu }%
,M\right) :\Omega \left( S\right) =0,~\forall S\in \Gamma \left( E,\pi
,M\right) \right\} .
\end{equation*}

The vector subbundle $\left( E^{0},\pi ^{0},M\right) $ will be called \emph{%
the annihilator vector subbundle of the IDS }$\left( E,\pi ,M\right) .$

\textbf{Proposition 3.2.1 }\emph{If }$\left( E,\pi ,M\right) $\emph{\ is an
IDS of the generalized Lie algebroid }%
\begin{equation*}
\left( \left( F,\nu ,N\right) ,\left[ ,\right] _{F,h},\left( \rho ,\eta
\right) \right)
\end{equation*}%
\emph{such that }$\Gamma \left( E,\pi ,M\right) =\left\langle
S_{1},...,S_{r}\right\rangle $\emph{, then it exists }$\Theta
^{r+1},...,\Theta ^{p}\in \Gamma \left( \overset{\ast }{h^{\ast }F},\overset{%
\ast }{h^{\ast }\nu },M\right) $\emph{\ linearly independent such that }$%
\Gamma \left( E^{0},\pi ^{0},M\right) =\left\langle \Theta ^{r+1},...,\Theta
^{p}\right\rangle .$

\textbf{Definition 3.2.2 }The IDS $\left( E,\pi ,M\right) $ of the
generalized Lie algebroid
\begin{equation*}
\left( \left( F,\nu ,N\right) ,\left[ ,\right] _{F,h},\left( \rho ,\eta
\right) \right)
\end{equation*}%
will be called \emph{involutive} if $\left[ S,T\right] _{h^{\ast }F}\in
\Gamma \left( E,\pi ,M\right) ,~$for any $S,T\in \Gamma \left( E,\pi
,M\right) .$

\textbf{Proposition 3.2.2 }\emph{If }$\left( E,\pi ,M\right) $\emph{\ is an
IDS of the generalized Lie algebroid }%
\begin{equation*}
\left( \left( F,\nu ,N\right) ,\left[ ,\right] _{F,h},\left( \rho ,\eta
\right) \right)
\end{equation*}%
\emph{and }$\left\{ S_{1},...,S_{r}\right\} $\emph{\ is a base for the }$%
\mathcal{F}\left( M\right) $\emph{-submodule }$\left( \Gamma \left( E,\pi
,M\right) ,+,\cdot \right) $\emph{\ then }$\left( E,\pi ,M\right) $\emph{\
is involutive if and only if }$\left[ S_{a},S_{b}\right] _{h^{\ast }F}\in
\Gamma \left( E,\pi ,M\right) ,~$for any $a,b\in \overline{1,r}.$

\section{Exterior differential calculus}

Let $\left( \left( F,\nu ,N\right) ,\left[ ,\right] _{F,h},\left( \rho ,\eta
\right) \right) $ be a generalized Lie algebroid. We denoted by $\Lambda
^{q}\left( F,\nu ,N\right) $ the set of \emph{differential forms of degree }$%
q.$ We remark that if
\begin{equation*}
\Lambda \left( F,\nu ,N\right) =\underset{q\geq 0}{\oplus }\Lambda
^{q}\left( F,\nu ,N\right) ,
\end{equation*}%
then we obtain \emph{the exterior differential algebra} $\left( \Lambda
\left( F,\nu ,N\right) ,+,\cdot ,\wedge \right) .$

\textbf{Definition 4.1 }For any $z\in \Gamma \left( F,\nu ,N\right) $, the
application
\begin{equation*}
\begin{array}{c}
\begin{array}{rcl}
\Lambda \left( F,\nu ,N\right) & ^{\underrightarrow{~\ \ L_{z}~\ \ }} &
\Lambda \left( F,\nu ,N\right)%
\end{array}%
,%
\end{array}%
\end{equation*}%
defined by%
\begin{equation*}
\begin{array}{c}
L_{z}\left( f\right) =\Gamma \left( Th\circ \rho ,h\circ \eta \right)
z\left( f\right) ,%
\end{array}%
\end{equation*}%
for any $f\in \mathcal{F}\left( N\right) $ and
\begin{equation*}
\begin{array}{cl}
L_{z}\omega \left( z_{1},...,z_{q}\right) & =\Gamma \left( Th\circ \rho
,h\circ \eta \right) z\left( \omega \left( \left( z_{1},...,z_{q}\right)
\right) \right) \\
& -\overset{q}{\underset{i=1}{\tsum }}\omega \left( \left( z_{1},...,\left[
z,z_{i}\right] _{F,h},...,z_{q}\right) \right) ,%
\end{array}%
\end{equation*}%
for any $\omega \in \Lambda ^{q}\mathbf{\ }\left( F,\nu ,N\right) $ and $%
z_{1},...,z_{q}\in \Gamma \left( F,\nu ,N\right) ,$ is called \emph{the
covariant Lie derivative with respect to the section }$z.$

\textbf{Theorem 4.1 }\emph{If }$z\in \Gamma \left( F,\nu ,N\right) ,$ $%
\omega \in \Lambda ^{q}\left( F,\nu ,N\right) $\emph{\ and }$\theta \in
\Lambda ^{r}\left( F,\nu ,N\right) $\emph{, then}%
\begin{equation*}
\begin{array}{c}
L_{z}\left( \omega \wedge \theta \right) =L_{z}\omega \wedge \theta +\omega
\wedge L_{z}\theta .%
\end{array}%
\leqno(4.1)
\end{equation*}

\textbf{Definition 4.2 }For any $z\in \Gamma \left( F,\nu ,N\right) $, the
application%
\begin{equation*}
\begin{array}{rcl}
\Lambda \left( F,\nu ,N\right) & ^{\underrightarrow{\ \ i_{z}\ \ }} &
\Lambda \left( F,\nu ,N\right) \\
\Lambda ^{q}\left( F,\nu ,N\right) \ni \omega & \longmapsto & i_{z}\omega
\in \Lambda ^{q-1}\left( F,\nu ,N\right) ,%
\end{array}%
\end{equation*}%
defined by $i_{z}f=0,$ for any $f\in \mathcal{F}\left( N\right) $ and
\begin{equation*}
\begin{array}{c}
i_{z}\omega \left( z_{2},...,z_{q}\right) =\omega \left(
z,z_{2},...,z_{q}\right) ,%
\end{array}%
\end{equation*}%
for any $z_{2},...,z_{q}\in \Gamma \left( F,\nu ,N\right) $, is called the
\emph{interior product associated to the section}~$z.$\bigskip

\textbf{Theorem 4.2 }\emph{If }$z\in \Gamma \left( F,\nu ,N\right) $\emph{,
then for any }$\omega \in $\emph{\ }$\Lambda ^{q}\left( F,\nu ,N\right) $%
\emph{\ and }$\theta \in $\emph{\ }$\Lambda ^{r}\left( F,\nu ,N\right) $%
\emph{\ we obtain}
\begin{equation*}
\begin{array}{c}
i_{z}\left( \omega \wedge \theta \right) =i_{z}\omega \wedge \theta +\left(
-1\right) ^{q}\omega \wedge i_{z}\theta .%
\end{array}%
\leqno(4.2)
\end{equation*}

\textbf{Theorem 4.3 }\emph{For any }$z,v\in \Gamma \left( F,\nu ,N\right) $%
\emph{\ we obtain}%
\begin{equation*}
\begin{array}{c}
L_{v}\circ i_{z}-i_{z}\circ L_{v}=i_{\left[ z,v\right] _{F,h}}.%
\end{array}%
\leqno(4.3)
\end{equation*}

\noindent

\textbf{Theorem 4.4 }\emph{The application }%
\begin{equation*}
\begin{array}{c}
\begin{array}{ccc}
\Lambda ^{q}\mathbf{\ }\left( F,\nu ,N\right)  & ^{\underrightarrow{\,\
d^{F}\,\ }} & \Lambda ^{q+1}\mathbf{\ }\left( F,\nu ,N\right)  \\
\omega  & \longmapsto  & d\omega
\end{array}%
\end{array}%
\end{equation*}%
\emph{defined by }$d^{F}f\left( z\right) =\Gamma \left( Th\circ \rho ,h\circ
\eta \right) \left( z\right) f,$ \emph{for any }$z\in \Gamma \left( F,\nu
,N\right) ,$ \emph{and}
\begin{equation*}
\begin{array}{l}
d^{F}\omega \left( z_{0},z_{1},...,z_{q}\right) =\overset{q}{\underset{i=0}{%
\tsum }}\left( -1\right) ^{i}\Gamma \left( Th\circ \rho ,h\circ \eta \right)
z_{i}\left( \omega \left( \left( z_{0},z_{1},...,\hat{z}_{i},...,z_{q}%
\right) \right) \right)  \\
~\ \ \ \ \ \ \ \ \ \ \ \ \ \ \ \ \ \ \ \ \ \ \ \ \ \ +\underset{i<j}{\tsum }%
\left( -1\right) ^{i+j}\omega \left( \left( \left[ z_{i},z_{j}\right]
_{F,h},z_{0},z_{1},...,\hat{z}_{i},...,\hat{z}_{j},...,z_{q}\right) \right) ,%
\end{array}%
\end{equation*}%
\emph{for any }$z_{0},z_{1},...,z_{q}\in \Gamma \left( F,\nu ,N\right) ,$
\emph{is unique with the following property:}%
\begin{equation*}
\begin{array}{c}
L_{z}=d^{F}\circ i_{z}+i_{z}\circ d^{F},~\forall z\in \Gamma \left( F,\nu
,N\right) .%
\end{array}%
\leqno(4.4)
\end{equation*}

This application will be called\emph{\ the exterior differentiation
ope\-ra\-tor for the exterior differential algebra of the generalized Lie
algebroid }$((F,\nu ,N),[,]_{F,h},(\rho ,\eta )).$

\textbf{Theorem 4.5} \emph{The exterior differentiation operator }$d^{F}$%
\emph{\ given by the previous theorem has the following properties:}\medskip

\noindent 1. \emph{\ For any }$\omega \in $\emph{\ }$\Lambda ^{q}\left(
F,\nu ,N\right) $\emph{\ and }$\theta \in $\emph{\ }$\Lambda ^{r}\left(
F,\nu ,N\right) $\emph{\ we obtain }%
\begin{equation*}
\begin{array}{c}
d^{F}\left( \omega \wedge \theta \right) =d^{F}\omega \wedge \theta +\left(
-1\right) ^{q}\omega \wedge d^{F}\theta .%
\end{array}%
\leqno(4.5)
\end{equation*}

\noindent 2.\emph{\ For any }$z\in \Gamma \left( F,\nu ,N\right) $ \emph{we
obtain }$L_{z}\circ d^{F}=d^{F}\circ L_{z}.$

\noindent3.\emph{\ }$d^{F}\circ d^{F}=0.$

\bigskip \textbf{Theorem 4.6} (of Maurer-Cartan type)\textbf{\ \ }

\emph{If }$((F,\nu ,N),[,]_{F,h},(\rho ,\eta ))$ \emph{is a generalized Lie
algebroid and} $d^{F}$\break \emph{is the ex\-te\-rior differentiation
operator for the exterior differential} $\mathcal{F}(N)$\emph{-algebra}%
\break $(\Lambda (F,\nu ,N),+,\cdot ,\wedge ),$ \emph{then we obtain the
structure equations of Maurer-Cartan type }%
\begin{equation*}
\begin{array}{c}
d^{F}t^{\alpha }=-\displaystyle\frac{1}{2}L_{\beta \gamma }^{\alpha
}t^{\beta }\wedge t^{\gamma },~\alpha \in \overline{1,p}%
\end{array}%
\leqno(\mathcal{C}_{1})
\end{equation*}%
\emph{and\ }%
\begin{equation*}
\begin{array}{c}
d^{F}\varkappa ^{\tilde{\imath}}=\theta _{\alpha }^{\tilde{\imath}}t^{\alpha
},~\tilde{\imath}\in \overline{1,n},%
\end{array}%
\leqno(\mathcal{C}_{2})
\end{equation*}%
\emph{where }$\left\{ t^{\alpha },\alpha \in \overline{1,p}\right\} ~$\emph{%
is the coframe of the vector bundle }$\left( F,\nu ,N\right) .$\bigskip
\noindent

This equations will be called \emph{the structure equations of Maurer-Cartan
type associa\-ted to the generalized Lie algebroid }$\left( \left( F,\nu
,N\right) ,\left[ ,\right] _{F,h},\left( \rho ,\eta \right) \right) .$

\bigskip\noindent\emph{Proof.} Let $\alpha \in \overline{1,p}$ be arbitrary.
Since
\begin{equation*}
\begin{array}{c}
d^{F}t^{\alpha }\left( t_{\beta },t_{\gamma }\right) =-L_{\beta \gamma
}^{\alpha },~\forall \beta ,\gamma \in \overline{1,p}%
\end{array}%
\end{equation*}%
it results that
\begin{equation*}
\begin{array}{c}
d^{F}t^{\alpha }=-\underset{\beta <\gamma }{\tsum }L_{\beta \gamma }^{\alpha
}t^{\beta }\wedge t^{\gamma }.%
\end{array}%
\leqno(1)
\end{equation*}
Since $L_{\beta \gamma }^{\alpha }=-L_{\gamma \beta }^{\alpha }$ and $%
t^{\beta }\wedge t^{\gamma }=-t^{\gamma }\wedge t^{\beta }$, for nay $\beta
,\gamma \in \overline{1,p},$ it results that
\begin{equation*}
\begin{array}{c}
\underset{\beta <\gamma }{\tsum }L_{\beta \gamma }^{\alpha }t^{\beta }\wedge
t^{\gamma }=\displaystyle\frac{1}{2}L_{\beta \gamma }^{\alpha }t^{\beta
}\wedge t^{\gamma }%
\end{array}%
\leqno(2)
\end{equation*}%
Using the equalities $\left( 1\right) $ and $\left( 2\right) $ it results
the structure equation $(\mathcal{C}_{1}).$

Let $\tilde{\imath}\in \overline{1,n}$ be arbitrarily. Since $d^{F}\varkappa
^{\tilde{\imath}}\left( t_{\alpha }\right) =\theta _{\alpha }^{\tilde{\imath}%
},~$\ for any $\alpha \in \overline{1,p},$ it results the structure equation
$(\mathcal{C}_{2}).$\hfill \emph{q.e.d.}\bigskip

\bigskip \noindent \textbf{Corollary 4.1 }\emph{If }$\left( \left( h^{\ast
}F,h^{\ast }\nu ,M\right) ,\left[ ,\right] _{h^{\ast }F},\left( \overset{%
h^{\ast }F}{\rho },Id_{M}\right) \right) $ \emph{\ is the pull-back Lie
algebroid associated to the generalized Lie algebroid} $\left( \left( F,\nu
,N\right) ,\left[ ,\right] _{F,h},\left( \rho ,\eta \right) \right) $ \emph{%
and }$d^{h^{\ast }F}$\emph{\ is the exterior differentiation operator for
the exterior differential }$\mathcal{F}\left( M\right) $\emph{-algebra}%
\break $\left( \Lambda \left( h^{\ast }F,h^{\ast }\nu ,M\right) ,+,\cdot
,\wedge \right) $\emph{, then we obtain the following structure equations of
Maurer-Cartan type \ }%
\begin{equation*}
\begin{array}{c}
d^{h^{\ast }F}T^{\alpha }=-\displaystyle\frac{1}{2}\left( L_{\beta \gamma
}^{\alpha }\circ h\right) T^{\beta }\wedge T^{\gamma },~\alpha \in \overline{%
1,p}%
\end{array}%
\leqno(\mathcal{C}_{1}^{\prime })
\end{equation*}%
\emph{and\ }%
\begin{equation*}
\begin{array}{c}
d^{h^{\ast }F}x^{i}=\left( \rho _{\alpha }^{i}\circ h\right) T^{\alpha
},~i\in \overline{1,m}.%
\end{array}%
\leqno(\mathcal{C}_{2}^{\prime })
\end{equation*}%
This equations will be called \emph{the structure equations of Maurer-Cartan
type associated to the pull-back Lie algebroid }%
\begin{equation*}
\left( \left( h^{\ast }F,h^{\ast }\nu ,M\right) ,\left[ ,\right] _{h^{\ast
}F},\left( \overset{h^{\ast }F}{\rho },Id_{M}\right) \right) .
\end{equation*}

\textbf{Theorem 4.7 (}of Cartan type)\textbf{\ }\emph{Let }$\left( E,\pi
,M\right) $\emph{\ be an IDS of the generalized Lie algebroid }$\left(
\left( F,\nu ,N\right) ,\left[ ,\right] _{F,h},\left( \rho ,\eta \right)
\right) .$\emph{\ If }$\left\{ \Theta ^{r+1},...,\Theta ^{p}\right\} $\emph{%
\ is a base for the }$\mathcal{F}\left( M\right) $\emph{-submodule }$\left(
\Gamma \left( E^{0},\pi ^{0},M\right) ,+,\cdot \right) $\emph{, then the IDS
}$\left( E,\pi ,M\right) $\emph{\ is involutive if and only if it exists }%
\begin{equation*}
\Omega _{\beta }^{\alpha }\in \Lambda ^{1}\left( h^{\ast }F,h^{\ast }\nu
,M\right) ,~\alpha ,\beta \in \overline{r+1,p}
\end{equation*}%
\emph{such that}
\begin{equation*}
d^{h^{\ast }F}\Theta ^{\alpha }=\Sigma _{\beta \in \overline{r+1,p}}\Omega
_{\beta }^{\alpha }\wedge \Theta ^{\beta }\in \mathcal{I}\left( \Gamma
\left( E^{0},\pi ^{0},M\right) \right) .
\end{equation*}

\emph{Proof. }Let $\left\{ S_{1},...,S_{r}\right\} $ be a base for the $%
\mathcal{F}\left( M\right) $-submodule $\left( \Gamma \left( E,\pi ,M\right)
,+,\cdot \right) $

Let $\left\{ S_{r+1},...,S_{p}\right\} \in \Gamma \left( h^{\ast }F,h^{\ast
}\nu ,M\right) $ such that $\left\{
S_{1},...,S_{r},S_{r+1},...,S_{p}\right\} $ is a base for the $\mathcal{F}%
\left( M\right) $-module
\begin{equation*}
\left( \Gamma \left( h^{\ast }F,h^{\ast }\nu ,M\right) ,+,\cdot \right) .
\end{equation*}

Let $\Theta ^{1},...,\Theta ^{r}\in \Gamma \left( \overset{\ast }{h^{\ast }F}%
,\overset{\ast }{h^{\ast }\nu },M\right) $ such that $\left\{ \Theta
^{1},...,\Theta ^{r},\Theta ^{r+1},...,\Theta ^{p}\right\} $ is a base for
the $\mathcal{F}\left( M\right) $-module
\begin{equation*}
\left( \Gamma \left( \overset{\ast }{h^{\ast }F},\overset{\ast }{h^{\ast
}\nu },M\right) ,+,\cdot \right) .
\end{equation*}

For any $a,b\in \overline{1,r}$ and $\alpha ,\beta \in \overline{r+1,p}$, we
have the equalities:%
\begin{equation*}
\begin{array}{ccc}
\Theta ^{a}\left( S_{b}\right) & = & \delta _{b}^{a} \\
\Theta ^{a}\left( S_{\beta }\right) & = & 0 \\
\Theta ^{\alpha }\left( S_{b}\right) & = & 0 \\
\Theta ^{\alpha }\left( S_{\beta }\right) & = & \delta _{\beta }^{\alpha }%
\end{array}%
\end{equation*}

We remark that the set of the $2$-forms%
\begin{equation*}
\left\{ \Theta ^{a}\wedge \Theta ^{b},\Theta ^{a}\wedge \Theta ^{\beta
},\Theta ^{\alpha }\wedge \Theta ^{\beta },~a,b\in \overline{1,r}\wedge
\alpha ,\beta \in \overline{r+1,p}\right\}
\end{equation*}%
is a base for the $\mathcal{F}\left( M\right) $-module $\left( \Lambda
^{2}\left( h^{\ast }F,h^{\ast }\nu ,M\right) ,+,\cdot \right) .$

Therefore, we have%
\begin{equation*}
d^{h^{\ast }F}\Theta ^{\alpha }=\Sigma _{b<c}A_{bc}^{\alpha }\Theta
^{b}\wedge \Theta ^{c}+\Sigma _{b,\gamma }B_{b\gamma }^{\alpha }\Theta
^{b}\wedge \Theta ^{\gamma }+\Sigma _{\beta <\gamma }C_{\beta \gamma
}^{\alpha }\Theta ^{\beta }\wedge \Theta ^{\gamma },\leqno\left( 1\right)
\end{equation*}%
where, $A_{bc}^{\alpha },B_{b\gamma }^{\alpha }$ and $C_{\beta \gamma
}^{\alpha },~a,b,c\in \overline{1,r},~\alpha ,\beta ,\gamma \in \overline{%
r+1,p}$ are real local functions such that $A_{bc}^{\alpha }=-A_{cb}^{\alpha
}$ and $C_{\beta \gamma }^{\alpha }=-C_{\gamma \beta }^{\alpha }.$

Using the formula%
\begin{equation*}
\begin{array}{cl}
d^{h^{\ast }F}\Theta ^{\alpha }\left( S_{b},S_{c}\right) & =\Gamma \left(
\overset{h^{\ast }F}{\rho },Id_{M}\right) S_{b}\left( \Theta ^{\alpha
}\left( S_{c}\right) \right) -\Gamma \left( \overset{h^{\ast }F}{\rho }%
,Id_{M}\right) S_{c}\left( \Theta ^{\alpha }\left( S_{b}\right) \right) \\
& -\Theta ^{\alpha }\left( \left[ S_{b},S_{c}\right] _{h^{\ast }F}\right) ,%
\end{array}%
\leqno\left( 2\right)
\end{equation*}%
we obtain that
\begin{equation*}
A_{bc}^{\alpha }=-\Theta ^{\alpha }\left( \left[ S_{b},S_{c}\right]
_{h^{\ast }F}\right) ,~\leqno\left( 3\right)
\end{equation*}%
for any $b,c\in \overline{1,r}$ and $\alpha \in \overline{r+1,p}.$

We admit that $\left( E,\pi ,M\right) $ is an involutive IDS of the
generalized Lie algebroid $\left( \left( F,\nu ,N\right) ,\left[ ,\right]
_{F,h},\left( \rho ,\eta \right) \right) .$

As $\left[ S_{b},S_{c}\right] _{h^{\ast }F}\in \Gamma \left( E,\pi ,M\right)
,$ for any $b,c\in \overline{1,r},$ it results that $\Theta ^{\alpha }\left( %
\left[ S_{b},S_{c}\right] _{h^{\ast }F}\right) =0,$ for any $b,c\in
\overline{1,r}$ and $\alpha \in \overline{r+1,p}.$ Therefore, for any $%
b,c\in \overline{1,r}$ and $\alpha \in \overline{r+1,p},$ we obtain $%
A_{bc}^{\alpha }=0$ and
\begin{equation*}
\begin{array}{ccl}
d^{h^{\ast }F}\Theta ^{\alpha } & = & \Sigma _{b,\gamma }B_{b\gamma
}^{\alpha }\Theta ^{b}\wedge \Theta ^{\gamma }+\frac{1}{2}C_{\beta \gamma
}^{\alpha }\Theta ^{\beta }\wedge \Theta ^{\gamma } \\
& = & \left( B_{b\gamma }^{\alpha }\Theta ^{b}+\frac{1}{2}C_{\beta \gamma
}^{\alpha }\Theta ^{\beta }\right) \wedge \Theta ^{\gamma }.%
\end{array}%
\end{equation*}

As
\begin{equation*}
\Omega _{\gamma }^{\alpha }\overset{put}{=}B_{b\gamma }^{\alpha }\Theta ^{b}+%
\frac{1}{2}C_{\beta \gamma }^{\alpha }\Theta ^{\beta }\in \Lambda ^{1}\left(
h^{\ast }F,h^{\ast }\nu ,M\right) ,~
\end{equation*}
for any $\alpha ,\beta \in \overline{r+1,p},$ it results the first
implication.

Conversely, we admit that it exists
\begin{equation*}
\Omega _{\beta }^{\alpha }\in \Lambda ^{1}\left( h^{\ast }F,h^{\ast }\nu
,M\right) ,~\alpha ,\beta \in \overline{r+1,p}
\end{equation*}%
such that
\begin{equation*}
d^{h^{\ast }F}\Theta ^{\alpha }=\Sigma _{\beta \in \overline{r+1,p}}\Omega
_{\beta }^{\alpha }\wedge \Theta ^{\beta },~\leqno\left( 4\right)
\end{equation*}%
for any $\alpha \in \overline{r+1,p}.$

Using the affirmations $\left( 1\right) ,\left( 2\right) $ and $\left(
4\right) $ we obtain that $A_{bc}^{\alpha }=0,~$for any $b,c\in \overline{1,r%
}$ and $\alpha \in \overline{r+1,p}.$

Using the affirmation $\left( 3\right) $, we obtain $\Theta ^{\alpha }\left( %
\left[ S_{b},S_{c}\right] _{h^{\ast }F}\right) =0,~$for any $b,c\in
\overline{1,r}$ and $\alpha \in \overline{r+1,p}.$

Therefore, we have $\left[ S_{b},S_{c}\right] _{h^{\ast }F}\in \Gamma \left(
E,\pi ,M\right) ,~$for any $b,c\in \overline{1,r}.$

Using the \emph{Proposition 3.2.2}, we obtain the second implication.\hfill
\emph{q.e.d.}\medskip

If $\left( \left( F^{\prime },\nu ^{\prime },N^{\prime }\right) ,\left[ ,%
\right] _{F^{\prime },h^{\prime }},\left( \rho ^{\prime },\eta ^{\prime
}\right) \right) $ is an another generalized Lie algebroid and $\left(
\varphi ,\varphi _{0}\right) $ is a $\mathbf{GLA}$-morphism of
\begin{equation*}
\left( \left( F,\nu ,N\right) ,\left[ ,\right] _{F,h},\left( \rho ,\eta
\right) \right)
\end{equation*}%
source and
\begin{equation*}
\left( \left( F^{\prime },\nu ^{\prime },N^{\prime }\right) ,\left[ ,\right]
_{F^{\prime },h^{\prime }},\left( \rho ^{\prime },\eta ^{\prime }\right)
\right)
\end{equation*}%
target, then obtain the application
\begin{equation*}
\begin{array}{ccc}
\Lambda ^{q}\left( F^{\prime },\nu ^{\prime },N^{\prime }\right) & ^{%
\underrightarrow{\ \left( \varphi ,\varphi _{0}\right) ^{\ast }\ }} &
\Lambda ^{q}\left( F,\nu ,N\right) \\
\omega ^{\prime } & \longmapsto & \left( \varphi ,\varphi _{0}\right) ^{\ast
}\omega ^{\prime }%
\end{array}%
,
\end{equation*}%
where
\begin{equation*}
\begin{array}{c}
\left( \left( \varphi ,\varphi _{0}\right) ^{\ast }\omega ^{\prime }\right)
\left( z_{1},...,z_{q}\right) =\omega ^{\prime }\left( \Gamma \left( \varphi
,\varphi _{0}\right) \left( z_{1}\right) ,...,\Gamma \left( \varphi ,\varphi
_{0}\right) \left( z_{q}\right) \right) ,%
\end{array}%
\end{equation*}%
for any $z_{1},...,z_{q}\in \Gamma \left( F,\nu ,N\right) .$

\textbf{Theorem 4.8 }\emph{If }$\left( \varphi ,\varphi _{0}\right) $\emph{\
is a }$\mathbf{GLA}$-\emph{morphism of }%
\begin{equation*}
\left( \left( F,\nu ,N\right) ,\left[ ,\right] _{F,h},\left( \rho ,\eta
\right) \right)
\end{equation*}%
\emph{\ source and }%
\begin{equation*}
\left( \left( F^{\prime },\nu ^{\prime },N^{\prime }\right) ,\left[ ,\right]
_{F^{\prime },h^{\prime }},\left( \rho ^{\prime },\eta ^{\prime }\right)
\right)
\end{equation*}%
\emph{target, then the following affirmations are satisfied:}\medskip

\noindent 1. \emph{For any }$\omega ^{\prime }\in \Lambda ^{q}\left(
F^{\prime },\nu ^{\prime },N^{\prime }\right) $\emph{\ and }$\theta ^{\prime
}\in \Lambda ^{r}\left( F^{\prime },\nu ^{\prime },N^{\prime }\right) $\emph{%
\ we obtain}%
\begin{equation*}
\begin{array}{c}
\left( \varphi ,\varphi _{0}\right) ^{\ast }\left( \omega ^{\prime }\wedge
\theta ^{\prime }\right) =\left( \varphi ,\varphi _{0}\right) ^{\ast }\omega
^{\prime }\wedge \left( \varphi ,\varphi _{0}\right) ^{\ast }\theta ^{\prime
}.%
\end{array}%
\leqno(4.13)
\end{equation*}

\noindent 2.\emph{\ For any }$z\in \Gamma \left( F,\nu ,N\right) $\emph{\
and }$\omega ^{\prime }\in \Lambda ^{q}\left( F^{\prime },\nu ^{\prime
},N^{\prime }\right) $\emph{\ we obtain}%
\begin{equation*}
\begin{array}{c}
i_{z}\left( \left( \varphi ,\varphi _{0}\right) ^{\ast }\omega ^{\prime
}\right) =\left( \varphi ,\varphi _{0}\right) ^{\ast }\left( i_{\varphi
\left( z\right) }\omega ^{\prime }\right) .%
\end{array}%
\leqno(4.14)
\end{equation*}

\noindent 3. \emph{If }$N=N^{\prime }$ \emph{and }$\left( Th\circ \rho
,h\circ \eta \right) =\left( Th^{\prime }\circ \rho ^{\prime },h^{\prime
}\circ \eta ^{\prime }\right) \circ \left( \varphi ,\varphi _{0}\right) ,$
\emph{then we obtain}
\begin{equation*}
\begin{array}{c}
\left( \varphi ,\varphi _{0}\right) ^{\ast }\circ d^{F^{\prime }}=d^{F}\circ
\left( \varphi ,\varphi _{0}\right) ^{\ast }.%
\end{array}%
\leqno(4.15)
\end{equation*}

\subsection{Exterior Differential Systems}

Let $\left( \left( h^{\ast }F,h^{\ast }\nu ,M\right) ,\left[ ,\right]
_{h^{\ast }F},\left( \overset{h^{\ast }F}{\rho },Id_{M}\right) \right) $ be
the pull-back Lie algebroid of the generalized Lie algebroid $\left( \left(
F,\nu ,N\right) ,\left[ ,\right] _{F,h},\left( \rho ,\eta \right) \right) $.

\textbf{Definition 4.1.1 }Any ideal $\left( \mathcal{I},+,\cdot \right) $ of
the exterior differential algebra of the pull-back Lie algebroid $\left(
\left( h^{\ast }F,h^{\ast }\nu ,M\right) ,\left[ ,\right] _{h^{\ast
}F},\left( \overset{h^{\ast }F}{\rho },Id_{M}\right) \right) $ closed under
differentiation operator $d^{h^{\ast }F},$ namely $d^{h^{\ast }F}\mathcal{%
I\subseteq I},$ will be called \emph{differential ideal of the generalized
Lie algebroid }$\left( \left( F,\nu ,N\right) ,\left[ ,\right] _{F,h},\left(
\rho ,\eta \right) \right) .$

In particular, if $h=Id_{N}=\eta $, then any ideal $\left( \mathcal{I}%
,+,\cdot \right) $ of the exterior differential algebra of the Lie algebroid
$\left( \left( F,\nu ,N\right) ,\left[ ,\right] _{F},\left( \rho
,Id_{M}\right) \right) $ closed under differentiation operator $d^{F}$ $,$
namely $d^{F}\mathcal{I\subseteq I},$ will be called \emph{differential
ideal of the Lie algebroid }$\left( \left( F,\nu ,N\right) ,\left[ ,\right]
_{F},\left( \rho ,Id_{M}\right) \right) .$

\textbf{Definition 4.1.2 }Let $\left( \mathcal{I},+,\cdot \right) $ be a
differential ideal of the generalized Lie algebroid $\left( \left( F,\nu
,N\right) ,\left[ ,\right] _{F,h},\left( \rho ,\eta \right) \right) $ or of
the Lie algebroid $\left( \left( F,\nu ,N\right) ,\left[ ,\right]
_{F},\left( \rho ,Id_{M}\right) \right) $ respectively$.$

If it exists an IDS $\left( E,\pi ,M\right) $ such that for all $k\in
\mathbb{N}^{\ast }$ and $\omega \in \mathcal{I}\cap \Lambda ^{k}\left(
h^{\ast }F,h^{\ast }\nu ,M\right) $ we have $\omega \left(
u_{1},...,u_{k}\right) =0,$ for any $u_{1},...,u_{k}\in \Gamma \left( E,\pi
,M\right) ,$ then we will say that $\left( \mathcal{I},+,\cdot \right) $%
\emph{\ is an exterior differential system (EDS) of the generalized Lie
algebroid }%
\begin{equation*}
\left( \left( F,\nu ,N\right) ,\left[ ,\right] _{F,h},\left( \rho ,\eta
\right) \right) .
\end{equation*}

In particular, if $h=Id_{N}=\eta $ and it exists an IDS $\left( E,\pi
,M\right) $ such that for all $k\in \mathbb{N}^{\ast }$ and $\omega \in
\mathcal{I}\cap \Lambda ^{k}\left( F,\nu ,M\right) $ we have $\omega \left(
u_{1},...,u_{k}\right) =0,$ for any $u_{1},...,u_{k}\in \Gamma \left( E,\pi
,M\right) ,$ then we will say that $\left( \mathcal{I},+,\cdot \right) $%
\emph{\ is an exterior differential system (EDS) of the Lie algebroid }$%
\left( \left( F,\nu ,N\right) ,\left[ ,\right] _{F},\left( \rho
,Id_{N}\right) \right) .$

\textbf{Theorem 4.1.1 (}of Cartan type) \emph{The IDS }$\left( E,\pi
,M\right) $\emph{\ of the generalized Lie algebroid }$\left( \left( F,\nu
,N\right) ,\left[ ,\right] _{F,h},\left( \rho ,\eta \right) \right) $\emph{\
is involutive, if and only if the ideal generated by the }$\mathcal{F}\left(
M\right) $\emph{-submodule }$\left( \Gamma \left( E^{0},\pi ^{0},M\right)
,+,\cdot \right) $\emph{\ is an EDS of the generalized Lie algebroid }$%
\left( \left( F,\nu ,N\right) ,\left[ ,\right] _{F,h},\left( \rho ,\eta
\right) \right) .$

\emph{Proof. }Let $\left( E,\pi ,M\right) $ be an involutive IDS of the
generalized Lie algebroid
\begin{equation*}
\left( \left( F,\nu ,N\right) ,\left[ ,\right] _{F,h},\left( \rho ,\eta
\right) \right) .
\end{equation*}

Let $\left\{ \Theta ^{r+1},...,\Theta ^{p}\right\} $ be a base for the $%
\mathcal{F}\left( M\right) $-submodule $\left( \Gamma \left( E^{0},\pi
^{0},M\right) ,+,\cdot \right) .$

We know that
\begin{equation*}
\mathcal{I}\left( \Gamma \left( E^{0},\pi ^{0},M\right) \right) =\cup _{q\in
\mathbb{N}}\left\{ \Omega _{\alpha }\wedge \Theta ^{\alpha },~\left\{ \Omega
_{r+1},...,\Omega _{p}\right\} \subset \Lambda ^{q}\left( h^{\ast }F,h^{\ast
}\nu ,M\right) \right\} .
\end{equation*}

Let $q\in \mathbb{N}$ and $\left\{ \Omega _{r+1},...,\Omega _{p}\right\}
\subset \Lambda ^{q}\left( h^{\ast }F,h^{\ast }\nu ,M\right) $ be arbitrary.

Using the \emph{Theorems 4.5 and 4.7} we obtain
\begin{equation*}
\begin{array}{ccl}
d^{h^{\ast }F}\left( \Omega _{\alpha }\wedge \Theta ^{\alpha }\right) & = &
d^{h^{\ast }F}\Omega _{\alpha }\wedge \Theta ^{\alpha }+\left( -1\right)
^{q+1}\Omega _{\beta }\wedge d^{h^{\ast }F}\Theta ^{\beta } \\
& = & \left( d^{h^{\ast }F}\Omega _{\alpha }+\left( -1\right) ^{q+1}\Omega
_{\beta }\wedge \Omega _{\alpha }^{\beta }\right) \wedge \Theta ^{\alpha }.%
\end{array}%
\end{equation*}

As
\begin{equation*}
d^{h^{\ast }F}\Omega _{\alpha }+\left( -1\right) ^{q+1}\Omega _{\beta
}\wedge \Omega _{\alpha }^{\beta }\in \Lambda ^{q+2}\left( h^{\ast
}F,h^{\ast }\nu ,M\right)
\end{equation*}%
it results that
\begin{equation*}
d^{h^{\ast }F}\left( \Omega _{\beta }\wedge \Theta ^{\beta }\right) \in
\mathcal{I}\left( \Gamma \left( E^{0},\pi ^{0},M\right) \right)
\end{equation*}

Therefore,
\begin{equation*}
d^{h^{\ast }F}\mathcal{I}\left( \Gamma \left( E^{0},\pi ^{0},M\right)
\right) \subseteq \mathcal{I}\left( \Gamma \left( E^{0},\pi ^{0},M\right)
\right) .
\end{equation*}

Conversely, let $\left( E,\pi ,M\right) $ be an IDS of the generalized Lie
algebroid
\begin{equation*}
\left( \left( F,\nu ,N\right) ,\left[ ,\right] _{F,h},\left( \rho ,\eta
\right) \right)
\end{equation*}%
such that the $\mathcal{F}\left( M\right) $-submodule $\left( \mathcal{I}%
\left( \Gamma \left( E^{0},\pi ^{0},M\right) \right) ,+,\cdot \right) $ is
an EDS of the generalized Lie algebroid $\left( \left( F,\nu ,N\right) ,%
\left[ ,\right] _{F,h},\left( \rho ,\eta \right) \right) .$

Let $\left\{ \Theta ^{r+1},...,\Theta ^{p}\right\} $ be a base for the $%
\mathcal{F}\left( M\right) $-submodule $\left( \Gamma \left( E^{0},\pi
^{0},M\right) ,+,\cdot \right) .$ As
\begin{equation*}
d^{h^{\ast }F}\mathcal{I}\left( \Gamma \left( E^{0},\pi ^{0},M\right)
\right) \subseteq \mathcal{I}\left( \Gamma \left( E^{0},\pi ^{0},M\right)
\right)
\end{equation*}%
it results that it exists
\begin{equation*}
\Omega _{\beta }^{\alpha }\in \Lambda ^{1}\left( h^{\ast }F,h^{\ast }\nu
,M\right) ,~\alpha ,\beta \in \overline{r+1,p}
\end{equation*}%
such that
\begin{equation*}
d^{h^{\ast }F}\Theta ^{\alpha }=\Sigma _{\beta \in \overline{r+1,p}}\Omega
_{\beta }^{\alpha }\wedge \Theta ^{\beta }\in \mathcal{I}\left( \Gamma
\left( E^{0},\pi ^{0},M\right) \right) .
\end{equation*}

Using the \emph{Theorem 4.7}, it results that $\left( E,\pi ,M\right) $ is
an involutive IDS.\hfill \emph{q.e.d.}\medskip

\addcontentsline{toc}{section}{References}

\begin{thebibliography}{9}
\bibitem{1} C. M. Arcus, \emph{The generalized Lie algebroids and their
applications}, arXiv: 1007.1451v2, math-ph, 10 Aug (2010).

\bibitem{2} R.L. Bryant, S.S. Chern, R.B. Gardner, H.L. Goldschmidt, and
P.A. Griffiths, \emph{Exterior Differential Systems,} Springer-Verlag, 1991.

\bibitem{3} J. Grabowski, P. Urbanski, \emph{Lie algebroids and
Poisson-Nijenhuis structures,} Rep. Math. Phys. \textbf{40} (1997), 196-208.

\bibitem{4} P Griffiths, \emph{Exterior Differential Systems and the
Calculus of Variations,} Progr. Math., No. 25, Birkh\"{a}user, Boston, MA,
1983.

\bibitem{5} P.J. Higgins, K. Mackenzei, \emph{Algebraic constructions in the
category of Lie algebroids}, J. Algebra, \textbf{129}, 194-230, (1990).

\bibitem{6} T. A. Ivey and J. M. Landsberg, \emph{Cartan for Beginners:
Differential Geometry via Moving Frames and Exterior Differential Systems,}
Graduate Texts in Mathematics, American Mathematical Society, 2003.

\bibitem{7} N. Kamran, \emph{An elementary introduction to exterior
differential systems,} In \textquotedblleft Geometric approaches to
differential equations (Canberra, 1995)\textquotedblright , volume 15 of
Austral. Math. Soc. Lect. Ser., Cambridge Univ. Press, Cambridge, (2000),
pages 100--115.

\bibitem{8} K. Mackenzie, \emph{Lie Grupoids and Lie Algebroids in
Differential Geometry,} London Math. Soc. Lecture Notes Series 124 Cambridge
Univ. Press, Cambridge, (1987).

\bibitem{9} C.M. Marle, \emph{Lie algebroids and Lie pseudoalgebras},
arxiv:math DG/0806.0919v2, (2008).
\end{thebibliography}

\bigskip

\bigskip

\hfill%
\begin{tabular}{c}
SECONDARY SCHOOL \textquotedblleft CORNELIUS RADU\textquotedblright, \\
RADINESTI VILLAGE, 217196, GORJ COUNTY, ROMANIA \\
e-mail: c\_arcus@yahoo.com, c\_arcus@radinesti.ro%
\end{tabular}

\end{document}
